% Source: https://artofproblemsolving.com/community/c3h1354536_flying_in_a_planein_a_cube
% Stable identity: aops-topic-1354536-q01
Put the cube as $[-1,1]^3$, centre $O$, and let the plane be $n\cdot p=0$ with $n=(n_1,n_2,n_3)$, $|n|=1$. The cross-section is a convex polygon, symmetric about $O$ (the cube and the plane both are), so it has an even number of sides; its sides lie in faces of the cube, one per face met, so it has $4$ or $6$ sides.

\emph{Squares.} A $4$-gon meets two opposite pairs of faces and misses (or only touches) the third pair, say $z=\pm1$. Then the section is the graph $z=ax+by$ over the square $[-1,1]^2$, with sides along $u=(2,0,2a)$ and $v=(0,2,2b)$. It is a square only if $u\cdot v=4ab=0$ and $|u|=|v|$, i.e.\ $a=b=0$: the plane is parallel to a face. This gives $3$ squares of area $4$.

\emph{Regular hexagons.} A hexagon meets all six faces; the side in the face $x=\pm1$ has direction $d_1=n\times e_1=(0,n_3,-n_2)$, and similarly $d_2=(-n_3,0,n_1)$, $d_3=(n_2,-n_1,0)$. In a regular hexagon the three side directions make angles of $60^\circ$ with each other, so for each pair $|\cos|=\dfrac{1}{2}$:
\[
4n_1^2n_2^2=\left(n_2^2+n_3^2\right)\left(n_1^2+n_3^2\right)\ \text{and cyclically}.
\]
With $A=n_1^2$, $B=n_2^2$, $C=n_3^2$ and $A+B+C=1$ these read $4AB=(1-A)(1-B)$, $4BC=(1-B)(1-C)$, $4CA=(1-C)(1-A)$. Dividing the first two gives $\dfrac{A}{1-A}=\dfrac{C}{1-C}$, so $A=C$; likewise $A=B$. So $A=B=C=\dfrac{1}{3}$ and $n\parallel(\pm1,\pm1,\pm1)$: the plane is perpendicular to one of the $4$ space diagonals. Such a plane, e.g.\ $x+y+z=0$ (in the figure, the cube is $[0,2]^3$ and the plane is $x-y+z=1$), passes through the midpoints of six edges, and the section is a regular hexagon with side $\sqrt2$ (distance between adjacent edge midpoints such as $(1,-1,0)$ and $(1,0,-1)$), of area $\dfrac{3\sqrt3}{2}\cdot2=3\sqrt3$.

\begin{asy}
usepackage("amsmath");
size(5cm);
defaultpen(fontsize(10pt));
pair pr(triple p){return (p.x+0.45*p.y, p.z+0.3*p.y);}
triple[] V={(0,0,0),(2,0,0),(2,2,0),(0,2,0),(0,0,2),(2,0,2),(2,2,2),(0,2,2)};
pair v(int i){return pr(V[i]);}
pair[] H={pr((0,0,1)),pr((1,0,0)),pr((2,1,0)),pr((2,2,1)),pr((1,2,2)),pr((0,1,2))};
fill(H[0]--H[1]--H[2]--H[3]--H[4]--H[5]--cycle,paleblue+opacity(0.6));
draw(v(0)--v(1)--v(2),linewidth(0.9)); draw(v(4)--v(5)--v(6)--v(7)--cycle,linewidth(0.9));
draw(v(0)--v(4),linewidth(0.9)); draw(v(1)--v(5),linewidth(0.9)); draw(v(2)--v(6),linewidth(0.9));
draw(v(0)--v(3),gray+dashed); draw(v(3)--v(2),gray+dashed); draw(v(3)--v(7),gray+dashed);
draw(H[0]--H[1],blue); draw(H[1]--H[2],blue+dashed); draw(H[2]--H[3],blue); draw(H[3]--H[4],blue+dashed); draw(H[4]--H[5],blue); draw(H[5]--H[0],blue+dashed);
for(int i=0;i<6;++i) dot(H[i],blue);
label("\mbox{}",(min(currentpicture,user=true)+max(currentpicture,user=true))/2,black); // keeps the text after the figure black when LaTeX typesets the labels inline
\end{asy}

The total area is $3\cdot4+4\cdot3\sqrt3=12\sqrt3+12$, so $a=12$, $b=3$, $c=12$ and $a+b+c=\boxed{27}$.
